% GENERATED FILE. Edit canonical lessons, then run npm run generate:books.
% canonical-source-hash: fnv1a64:acbfd8e06d46cb18
% canonical-lessons: KA-C480-bande, KA-C480-marubhumi, KA-C480-bayalu, KA-C480-prasthabhumi, KA-C480-save

\chapter{Rocks, Deserts and Plains}
\label{ch:ka-rocks-deserts-and-plains}
\hlchaptermodality{480}

\begin{chapteropening}
\textbf{By the end of this chapter:} \emph{I can say a rock, a desert, an open plain, a plateau and to wear away.}

Complete the last lesson of chapter 480: I can say a rock, a desert, an open plain, a plateau and to wear away.
\end{chapteropening}

\section[\texorpdfstring{\kn{ಬಂಡೆ} (baṇḍe)}{ಬಂಡೆ (baṇḍe)}]{\kn{ಬಂಡೆ} (baṇḍe) — a rock; a boulder}

\label{lesson:KA-C480-bande}



\begin{quote}
Before the new one: say the Kannada for to sprout, then the Kannada for a swamp.
\end{quote}

\subsection*{You'll want to know: \kn{ಬಂಡೆ}}
\textbf{\kn{ಬಂಡೆ}} — \emph{baṇḍe} — \textquotedblleft{}a rock; a boulder\textquotedblright{}.

\begin{quote}
\kn{ನದಿಯ} \kn{ನಡುವೆ} \kn{ಒಂದು} \kn{ದೊಡ್ಡ} \kn{ಬಂಡೆ} \kn{ಇದೆ}.
\end{quote}

\begin{quote}
\emph{nadiya naḍuve ondu doḍḍa baṇḍe ide.}
\end{quote}

\begin{quote}
There is a big rock in the middle of the river.
\end{quote}

Plural: \kn{ಬಂಡೆಗಳು}, \emph{baṇḍegaḷu}.

\kn{ಕಲ್ಲು} is a stone; \kn{ಬಂಡೆ} is a big rock or a slab.

\begin{cousinweb}
\textbf{\kn{ಬಂಡೆ}} (\emph{baṇḍe}) is native Dravidian: DEDR 3903 lists Kannada \emph{baṇḍe}, \textquotedblleft{}a rock, a block or slab of stone\textquotedblright{}.
\end{cousinweb}

\begin{grammarlens}[title={What you've built}]
One more word for getting around.
\end{grammarlens}

\subsection*{Your turn}
\begin{itemize}
\raggedright
  \item \emph{Say it:} \emph{baṇḍe}
  \item \emph{Say it:} \emph{baṇḍe}, once more
  \item \emph{Say it:} \emph{nadiya naḍuve ondu doḍḍa baṇḍe ide.} — There is a big rock in the middle of the river.
  \item \emph{Recall:} say the Kannada for to sprout, then the Kannada for a swamp, then say \emph{baṇḍe} again
\end{itemize}

\begin{tcolorbox}[breakable,colback=teal!4,colframe=teal!35!black,title={Before you move on}]
What does \kn{ಬಂಡೆ} mean? (\textquotedblleft{}A rock; a boulder\textquotedblright{}.) Say it once more.
\end{tcolorbox}
\section[\texorpdfstring{\kn{ಮರುಭೂಮಿ} (marubhūmi)}{ಮರುಭೂಮಿ (marubhūmi)}]{\kn{ಮರುಭೂಮಿ} (marubhūmi) — a desert}

\label{lesson:KA-C480-marubhumi}



\begin{quote}
Before the new one: say the Kannada for a swamp, then the Kannada for a rock.
\end{quote}

\subsection*{You'll want to know: \kn{ಮರುಭೂಮಿ}}
\textbf{\kn{ಮರುಭೂಮಿ}} — \emph{marubhūmi} — \textquotedblleft{}a desert\textquotedblright{}.

\begin{quote}
\kn{ಮರುಭೂಮಿಯಲ್ಲಿ} \kn{ಒಂಟೆಗಳು} \kn{ನೀರಿಲ್ಲದೆ} \kn{ತುಂಬಾ} \kn{ದಿನ} \kn{ಇರುತ್ತವೆ}.
\end{quote}

\begin{quote}
\emph{marubhūmiyalli oṇṭegaḷu nīrillade tumbā dina iruttave.}
\end{quote}

\begin{quote}
In the desert camels go many days without water.
\end{quote}

\kn{ಮರು}, a wilderness, plus \kn{ಭೂಮಿ}, land.

\begin{cousinweb}
\textbf{\kn{ಮರುಭೂಮಿ}} (\emph{marubhūmi}) is built from two Sanskrit words: \emph{maru}, \textquotedblleft{}a wilderness\textquotedblright{} (CDIAL 9876a), and \emph{bhūmi}, \textquotedblleft{}earth, ground\textquotedblright{} (CDIAL 9557).
\end{cousinweb}

\begin{grammarlens}[title={What you've built}]
One more word for getting around.
\end{grammarlens}

\subsection*{Your turn}
\begin{itemize}
\raggedright
  \item \emph{Say it:} \emph{marubhūmi}
  \item \emph{Say it:} \emph{marubhūmi}, once more
  \item \emph{Say it:} \emph{marubhūmiyalli oṇṭegaḷu nīrillade tumbā dina iruttave.} — In the desert camels go many days without water.
  \item \emph{Recall:} say the Kannada for a swamp, then the Kannada for a rock, then say \emph{marubhūmi} again
\end{itemize}

\begin{tcolorbox}[breakable,colback=teal!4,colframe=teal!35!black,title={Before you move on}]
What does \kn{ಮರುಭೂಮಿ} mean? (\textquotedblleft{}A desert\textquotedblright{}.) Say it once more.
\end{tcolorbox}
\section[\texorpdfstring{\kn{ಬಯಲು} (bayalu)}{ಬಯಲು (bayalu)}]{\kn{ಬಯಲು} (bayalu) — an open plain; open ground}

\label{lesson:KA-C480-bayalu}



\begin{quote}
Before the new one: say the Kannada for a rock, then the Kannada for a desert.
\end{quote}

\subsection*{You'll want to know: \kn{ಬಯಲು}}
\textbf{\kn{ಬಯಲು}} — \emph{bayalu} — \textquotedblleft{}an open plain; open ground\textquotedblright{}.

\begin{quote}
\kn{ಊರಿನ} \kn{ಹೊರಗೆ} \kn{ದೊಡ್ಡ} \kn{ಬಯಲು} \kn{ಇದೆ}.
\end{quote}

\begin{quote}
\emph{ūrina horage doḍḍa bayalu ide.}
\end{quote}

\begin{quote}
Outside the town there is a big open plain.
\end{quote}

\kn{ಬಯಲುಸೀಮೆ}, the plains country, as against the hills.

\begin{cousinweb}
\textbf{\kn{ಬಯಲು}} (\emph{bayalu}) is native Dravidian: DEDR 5258 lists Kannada \emph{bayal}, \emph{bayil}, \textquotedblleft{}open space, a field, a plain\textquotedblright{}.
\end{cousinweb}

\begin{grammarlens}[title={What you've built}]
One more word for getting around.
\end{grammarlens}

\subsection*{Your turn}
\begin{itemize}
\raggedright
  \item \emph{Say it:} \emph{bayalu}
  \item \emph{Say it:} \emph{bayalu}, once more
  \item \emph{Say it:} \emph{ūrina horage doḍḍa bayalu ide.} — Outside the town there is a big open plain.
  \item \emph{Recall:} say the Kannada for a rock, then the Kannada for a desert, then say \emph{bayalu} again
\end{itemize}

\begin{tcolorbox}[breakable,colback=teal!4,colframe=teal!35!black,title={Before you move on}]
What does \kn{ಬಯಲು} mean? (\textquotedblleft{}An open plain; open ground\textquotedblright{}.) Say it once more.
\end{tcolorbox}
\section[\texorpdfstring{\kn{ಪ್ರಸ್ಥಭೂಮಿ} (prasthabhūmi)}{ಪ್ರಸ್ಥಭೂಮಿ (prasthabhūmi)}]{\kn{ಪ್ರಸ್ಥಭೂಮಿ} (prasthabhūmi) — a plateau; a tableland}

\label{lesson:KA-C480-prasthabhumi}



\begin{quote}
Before the new one: say the Kannada for a desert, then the Kannada for an open plain.
\end{quote}

\subsection*{You'll want to know: \kn{ಪ್ರಸ್ಥಭೂಮಿ}}
\textbf{\kn{ಪ್ರಸ್ಥಭೂಮಿ}} — \emph{prasthabhūmi} — \textquotedblleft{}a plateau; a tableland\textquotedblright{}.

\begin{quote}
\kn{ಈ} \kn{ಪ್ರಸ್ಥಭೂಮಿಯಲ್ಲಿ} \kn{ಮಳೆ} \kn{ಕಡಿಮೆ}.
\end{quote}

\begin{quote}
\emph{ī prasthabhūmiyalli maḷe kaḍime.}
\end{quote}

\begin{quote}
There is little rain on this plateau.
\end{quote}

\kn{ಪ್ರಸ್ಥ}, flat land, plus \kn{ಭೂಮಿ}, land.

\begin{cousinweb}
\textbf{\kn{ಪ್ರಸ್ಥಭೂಮಿ}} (\emph{prasthabhūmi}) is a modern compound built from Sanskrit words: \emph{prastha}, \textquotedblleft{}flat land\textquotedblright{} (CDIAL 8868), and \emph{bhūmi}, \textquotedblleft{}earth, land\textquotedblright{} (CDIAL 9557).
\end{cousinweb}

\begin{grammarlens}[title={What you've built}]
One more word for getting around.
\end{grammarlens}

\subsection*{Your turn}
\begin{itemize}
\raggedright
  \item \emph{Say it:} \emph{prasthabhūmi}
  \item \emph{Say it:} \emph{prasthabhūmi}, once more
  \item \emph{Say it:} \emph{ī prasthabhūmiyalli maḷe kaḍime.} — There is little rain on this plateau.
  \item \emph{Recall:} say the Kannada for a desert, then the Kannada for an open plain, then say \emph{prasthabhūmi} again
\end{itemize}

\begin{tcolorbox}[breakable,colback=teal!4,colframe=teal!35!black,title={Before you move on}]
What does \kn{ಪ್ರಸ್ಥಭೂಮಿ} mean? (\textquotedblleft{}A plateau; a tableland\textquotedblright{}.) Say it once more.
\end{tcolorbox}
\section[\texorpdfstring{\kn{ಸವೆ} (save)}{ಸವೆ (save)}]{\kn{ಸವೆ} (save) — to wear away; to erode}

\label{lesson:KA-C480-save}



\begin{quote}
Before the new one: say the Kannada for an open plain, then the Kannada for a plateau.
\end{quote}

\subsection*{You'll want to know: \kn{ಸವೆ}}
\textbf{\kn{ಸವೆ}} — \emph{save} — \textquotedblleft{}to wear away; to erode\textquotedblright{}.

\begin{quote}
\kn{ತುಂಬಾ} \kn{ನಡೆದಿದ್ದರಿಂದ} \kn{ನನ್ನ} \kn{ಚಪ್ಪಲಿ} \kn{ಸವೆಯಿತು}.
\end{quote}

\begin{quote}
\emph{tumbā naḍediddarinda nanna cappali saveyitu.}
\end{quote}

\begin{quote}
From so much walking my sandals wore out.
\end{quote}

Past (it …): \kn{ಸವೆಯಿತು}, \emph{saveyitu}.

\kn{ಸವೆತ}, wear or erosion: \kn{ಮಣ್ಣಿನ} \kn{ಸವೆತ}, soil erosion.

\begin{cousinweb}
\textbf{\kn{ಸವೆ}} (\emph{save}) is native Dravidian: DEDR 2343 lists Kannada \emph{same}, \emph{save}, \textquotedblleft{}to be consumed, to wear away, to be abraded\textquotedblright{}.
\end{cousinweb}

\begin{grammarlens}[title={What you've built}]
One more word for getting around.
\end{grammarlens}

\subsection*{Your turn}
\begin{itemize}
\raggedright
  \item \emph{Say it:} \emph{save}
  \item \emph{Say it:} \emph{save}, once more
  \item \emph{Say it:} \emph{tumbā naḍediddarinda nanna cappali saveyitu.} — From so much walking my sandals wore out.
  \item \emph{Recall:} say the Kannada for an open plain, then the Kannada for a plateau, then say \emph{save} again
\end{itemize}

\begin{tcolorbox}[breakable,colback=teal!4,colframe=teal!35!black,title={Before you move on}]
What does \kn{ಸವೆ} mean? (\textquotedblleft{}To wear away; to erode\textquotedblright{}.) Say it once more.
\end{tcolorbox}
